\documentclass[11pt]{article}

\usepackage[margin=1in]{geometry}
\usepackage{amsmath,amssymb,amsthm,mathtools}
\usepackage{microtype}
\usepackage{pgfplots}
\pgfplotsset{compat=1.17}

\newtheorem{theorem}{Theorem}
\newtheorem{lemma}{Lemma}
\newtheorem{remark}{Remark}

\newcommand{\LS}{\mathrm{LS}}
\newcommand{\OPT}{\mathrm{OPT}}

\title{A m dependant Upper Bound}
\author{}
\date{}

\begin{document}
\maketitle

\section{Assumptions}

We only consider \(m\ge 3\) identical parallel machines, as we have previously shown, 
that we have found the best possible upper bound for $m=2$. and jobs
 \(M_i(S)\) be the set of jobs assigned to machine
\(i\).  If \(j\in M_i(S)\), define
\[
 C_j^S(\alpha)
 =
 \sum_{\substack{k\in M_i(S)\\k\text{ before }j}}p_k+\alpha p_j,
 \qquad 0<\alpha\le\frac12.
\]


We will make use of the following results.
\begin{enumerate}
 \item Any $\lambda$ maximizing instance $I$ has the property 
 that every machine contains at most \(m\) jobs in both $S_{WORST}$ and $S_{OPT}$.

 \item There exists a $\lambda$ maximizing instance $I$ such that 
 Some job \(p\) is single scheduled  in $S_{OPT}$, and there exists a list such that \(p\) is the
 last job scheduled in $S_{WORST}$.
 Additionally, no job in $I$ finishes strictly after $p$ in $S_{WORST}$.

 \item $\max_{I \in P||\sum_{j \in I}p_jC_j(\alpha)}\lambda(I)(\alpha) = \frac43$. 
 \item We assume $n > m$. Else $\lambda(I)(\alpha) =1$ and we are done. 

\end{enumerate}

Next, we derive a useful form of the objective.
Fix a machine $i$ and write its jobs in order of processing in schedule $S$:
\[
 j_{i,1},j_{i,2},\ldots,j_{i,r_i}.
\]
Their incurred cost is: 
\[
\begin{aligned}
\sum_{s=1}^{r_i}p_{j_{i,s}}C_{j_{i,s}}^S(\alpha)
&=
\sum_{s=1}^{r_i}
p_{j_{i,s}}
\left(
\sum_{t=1}^{s-1}p_{j_{i,t}}
+\alpha p_{j_{i,s}}
\right)\\
&=
\alpha\sum_{s=1}^{r_i}p_{j_{i,s}}^2
+
\sum_{s=1}^{r_i}\sum_{t=1}^{s-1}
p_{j_{i,s}}p_{j_{i,t}}\\
&=
\alpha\sum_{j\in M_i(S)}p_j^2
+
\sum_{\{j,k\}\subseteq M_i(S)}p_jp_k.
\end{aligned}
\]

This derivation stems again from Schwiegelshohn (hier reference rein)
Summing this identity over all machines gives
\[
 C(S)(\alpha)=\alpha Q+P(S).
\tag{1}
\]

where: 
\[
 Q:=\sum_jp_j^2
\]
and, for every schedule \(S\),
\[
 P(S):=
 \sum_{i=1}^m\ \sum_{\{j,k\}\subseteq M_i(S)}p_jp_k.
\]

 We write
\[
 P_{W}:=P(S_{WORST}),\qquad P^*:=P(S_{OPT}).
\]

Because \(\OPT\) has one single scheduled job and every machine
contains at most \(m\) jobs,
\[
 n\le 1+m(m-1).
\tag{2}
\]

\section{Max. Cost difference}

Define
\[
 \mu_m:=\sqrt{\frac{m-1}{2}},
 \qquad
 K_m:=(m-1)^2,
 \qquad
 C_m:=m\left(1+(K_m-1)\mu_m\right).
\]

\begin{lemma}
\label{lemma: P_W-P*}
Under the assumptions above,
\[
 P_{W}-P^*\le C_m\sqrt{QP^*}.
\tag{3}
\]
\end{lemma}

\begin{proof}
Consider the jobs in list order of $S_W$.  When job \(j\) is scheduled, let \(\ell_j\)
be the current load of the receiving machine (eq. this is the starting time of j).  The placement creates pair
cost \(p_j\ell_j\), and therefore
\[
 P_{W}=\sum_jp_j\ell_j.
\]

Fix any job \(j\) with \(\ell_j>0\) (this exists as we schedule at least m+1 jobs). 
We consider now only the partial schedule consisting of all jobs scheduled before 
(and including) $j$.
Since $j$ is assigned to a globally least-loaded machine,
every machine has current load at least \(\ell_j\).
Choose a largest currently assigned job from each machine (excluding j).  At that moment
each machine contains at most \(m\) jobs, so every selected job has size at
least
\[
 \frac{\ell_j}{m}.
\]
The \(m\) selected jobs and \(j\) are \(m+1\) distinct jobs.  Thus, b.c of pidgeonhole, 
two of them share a machine in $S^*$.

If the colliding pair consists of two selected jobs, then
\[
 P^*\ge\frac{\ell_j^2}{m^2}.
\]
Since \(p_j\le\sqrt Q\), it follows that
\[
 p_j\ell_j\le m\sqrt{QP^*}.
\tag{4}
\]

If the colliding pair consists of \(j\) and a selected job, then
\[
 P^*\ge\frac{p_j\ell_j}{m},
\]
and hence
\[
 p_j\ell_j\le mP^*.
\tag{5}
\]

For any \(k\) positive numbers \(x_1,\ldots,x_k\),
\[
 \sum_{r<s}x_rx_s
 \le \frac{k-1}{2}\sum_{r=1}^kx_r^2.
\tag{6}
\]
Indeed, this follows from
\((\sum_rx_r)^2\le k\sum_rx_r^2\).

(ref: J. Michael Steele, The Cauchy–Schwarz Master Class: An Introduction to the Art of Mathematical Inequalities, Cambridge University Press, 2004)
Applying (6) to every  set $M_{i}(S_{OPT})$ and using $|M_{i}(S_{OPT})|\le m$, then summing over all machines, gives: 
\[
 P^*\le\frac{m-1}{2}Q.
\tag{7}
\]
From (5), we have
\[
p_j\ell_j \le mP^*.
\]

Moreover, (7) gives
\[
P^* \le \frac{m-1}{2}Q.
\]

Multiplying (7) by \(P^*\ge 0\), we obtain
\[
(P^*)^2
\le
\frac{m-1}{2}QP^*.
\]

Taking square roots on both sides yields
\[
P^*
\le
\sqrt{\frac{m-1}{2}}\sqrt{QP^*}.
\]

Substituting this estimate into (5), we conclude that
\[
\begin{aligned}
p_j\ell_j
&\le mP^*\\
&\le
m\sqrt{\frac{m-1}{2}}\sqrt{QP^*}.
\end{aligned}
\tag{8}
\]

Although direct substitution of (7) into (5) also gives the valid bound
\[
p_j\ell_j\le \frac{m(m-1)}{2}Q,
\]
the mixed estimate in (8) is more useful because it retains the dependence
on \(P^*\). In particular, with \(z=P^*/Q\),
\[
\frac{\sqrt{QP^*}}{\alpha Q+P^*}
=
\frac{\sqrt z}{\alpha+z}.
\]

Since
\[
\max_{z\ge0}\frac{\sqrt z}{\alpha+z}
=
\frac{1}{2\sqrt{\alpha}},
\]
the mixed bound leads to an approximation estimate of order
\(1/\sqrt{\alpha}\), whereas a bound involving only \(Q\) would remain of
order \(1/\alpha\).

Consequently, every placement with \(\ell_j>0\) satisfies
\[
 p_j\ell_j\le m\mu_m\sqrt{QP^*}.
\tag{9}
\]

The first \(m\) jobs of an ordinary list schedule are placed on empty
machines and have \(\ell_j=0\).  By (2), the number of placements with strictly
positive $\ell_j$ is at most
\[
 n-m\le 1+(m-1)m - m = (m-1)^2=K_m.
\]

We now use both special properties of \(p\). Since \(p\) is last in some list of $S_{W}$,
and at least $m$ jobs precede it in the schedule. its preceding load \(\ell_p\) is positive.  Immediately before placing \(p\),
choose a largest job from each of the \(m\) machines in $S_W$.  None of these
selected jobs can lie on the singleton machine of $S^*$, because that
machine contains only \(p\).  Hence the \(m\) selected jobs occupy only
\(m-1\) optimal machines, and again with pidgeonhole, two selected jobs share a machine in $S^*$.
Therefore, (4) applies to \(p\), giving: 
estimate
\[
 p\ell_p\le m\sqrt{QP^*}.
\tag{10}
\]

There are at most \(K_m-1\) other positive-load ($\ell_j > 0$) assignments .  Summing (9)
over those placements and adding (10), we obtain
\[
\begin{aligned}
 P_{W}
 &\le
 m\left(1+(K_m-1)\mu_m\right)\sqrt{QP^*}\\
 &=C_m\sqrt{QP^*}.
\end{aligned}
\]
Finally \(P_{\LS}-P^*\le P_{\LS}\), proving (3).
\end{proof}

\section{The upper bound}

Expanding the established result for $\alpha = \frac12$ by means of (1) gives
\[
 \frac12Q+P_{\LS}
 \le
 \frac43\left(\frac12Q+P^*\right),
\]
and hence
\[
 P_{\LS}-P^*\le\frac16Q+\frac13P^*.
\tag{11}
\]
We now translate these two bounds on the additive cost difference into
bounds on the approximation ratio. By (1),
\[
\begin{aligned}
\frac{C(S_{WORST})(\alpha)}{C(S_{OPT})(\alpha)}
&=
\frac{\alpha Q+P_W}{\alpha Q+P^*}\\
&=
1+\frac{P_W-P^*}{\alpha Q+P^*}.
\end{aligned}
\]
Using (11) to bound the numerator gives
\[
\frac{C(S_{WORST})(\alpha)}{C(S_{OPT})(\alpha)}
\le
1+
\frac{\frac16Q+\frac13P^*}{\alpha Q+P^*}.
\]
Using Lemma~\ref{lemma: P_W-P*} instead gives
\[
\frac{C(S_{WORST})(\alpha)}{C(S_{OPT})(\alpha)}
\le
1+
\frac{C_m\sqrt{QP^*}}{\alpha Q+P^*}.
\]
Both inequalities hold simultaneously. We can therefore take the smaller
of the two right-hand sides:
\[
\frac{C(S_{WORST})(\alpha)}{C(S_{OPT})(\alpha)}
\le
1+
\min\left\{
\frac{\frac16Q+\frac13P^*}{\alpha Q+P^*},
\frac{C_m\sqrt{QP^*}}{\alpha Q+P^*}
\right\}.
\]
Finally, since \(Q>0\), define
\[
 z:=\frac{P^*}{Q}.
\]
Dividing the numerator and denominator of both fractions by \(Q\), and
using
\[
\frac{\sqrt{QP^*}}{Q}
=
\sqrt{\frac{P^*}{Q}}
=\sqrt z,
\]
gives
\[
 \frac{C(S_{WORST})(\alpha)}{C(S_{OPT})(\alpha)}
 \le
 1+
 \min\left\{
 \frac{\frac16+\frac13z}{\alpha+z},
 \frac{C_m\sqrt z}{\alpha+z}
 \right\}.
\tag{12}
\]
Next we try to determine, for which values of $z$, which of the above expressions is minimal. 
TO this end, define \(z_m\)  to be the smaller positive solution of
\[
 \frac16+\frac13z=C_m\sqrt z.
\tag{13}
\]. 

Equivalently, define
\[
 y_m:=
 \frac{1}{3C_m+\sqrt{9C_m^2-2}},
 \qquad
 z_m:=y_m^2.
\tag{14}
\]

Note that, since $C_m > 1$, $z_m$ is indeed real and positive. 

\begin{theorem}
\label{thm:main}
For every \(0<\alpha\le1/2\),
\[
 \lambda(I)(\alpha)
 \le U_m(\alpha),
\tag{15}
\]
where
\[
 U_m(\alpha)=
 \begin{cases}
 \displaystyle
 1+\frac{C_m}{2\sqrt\alpha},
 &0<\alpha\le z_m,\\[3mm]
 \displaystyle
 1+\frac{C_m\sqrt{z_m}}{\alpha+z_m}
 =
 1+\frac{1+2z_m}{6(\alpha+z_m)},
 &z_m\le\alpha\le\frac12.
 \end{cases}
\tag{16}
\]
Moreover,
\[
 U_m(\alpha)<1+\frac{1}{6\alpha}
 \qquad\text{for }0<\alpha<\frac12,
\tag{17}
\]
and
\[
 U_m\left(\frac12\right)=\frac43.
\tag{18}
\]
\end{theorem}

\begin{proof}
For fixed \(\alpha\), define
\[
 g_1(z):=\frac{\frac16+\frac13z}{\alpha+z},
 \qquad
 g_2(z):=\frac{C_m\sqrt z}{\alpha+z}.
\]
For \(0<\alpha\le1/2\),
\[
 g_1'(z)=\frac{\alpha/3-1/6}{(\alpha+z)^2}\le0.
\]
By (13), \(g_1(z_m)=g_2(z_m)\).

Suppose first that \(\alpha\ge z_m\).  The function \(g_2\) is increasing
on \([0,\alpha]\), and hence on \([0,z_m]\).  Thus, for \(z\le z_m\),
\[
 \min\{g_1(z),g_2(z)\}\le g_2(z_m),
\]
while for \(z\ge z_m\), monotonicity of \(g_1\) gives
\[
 \min\{g_1(z),g_2(z)\}\le g_1(z)\le g_1(z_m).
\]
This proves the second branch of (16).

Suppose next that \(0<\alpha\le z_m\).  The global maximum of \(g_2\) occurs
at \(z=\alpha\), where
\[
 g_2(\alpha)=\frac{C_m}{2\sqrt\alpha}.
\]
For \(z\le z_m\), the minimum in (12) is at most this global maximum.  For
\(z\ge z_m\), it is at most
\[
 g_1(z_m)=g_2(z_m)\le g_2(\alpha).
\]
This proves the first branch of (16), and therefore (15).

\begin{figure}[ht]
\centering
\begin{minipage}{0.48\textwidth}
\centering
\begin{tikzpicture}
% Numerical values used only to render the qualitative plot.
% The visible labels remain symbolic.
\def\alphaval{0.08}
\def\Cmplot{1.0}
\def\zone{0.03136}
\begin{semilogxaxis}[
    width=\textwidth,
    height=0.72\textwidth,
    axis lines=left,
    xlabel={},
    ylabel={value},
    xmin=0.008,
    xmax=20,
    ymin=0,
    ymax=3.0,
    samples=400,
    domain=0.008:20,
    legend style={draw=none, fill=none, at={(0.98,0.98)}, anchor=north east, font=\scriptsize},
    xtick={\zone,\alphaval},
    xticklabels={$z_m$,$\alpha$},
    ytick=\empty,
    tick style={draw=none},
    clip=false,
]
  \addplot[black, thick]
    {(1/6 + x/3)/(\alphaval + x)};
  \addlegendentry{$g_1(z)=\frac{1/6+z/3}{\alpha+z}$}

  \addplot[black, thick, dashed]
    {\Cmplot*sqrt(x)/(\alphaval + x)};
  \addlegendentry{$g_2(z)=\frac{C_m\sqrt z}{\alpha+z}$}

  \draw[densely dotted] (axis cs:\zone,0) --
    (axis cs:\zone,{(1/6+\zone/3)/(\alphaval+\zone)});
  \draw[densely dotted] (axis cs:\alphaval,0) --
    (axis cs:\alphaval,{\Cmplot*sqrt(\alphaval)/(\alphaval+\alphaval)});
\end{semilogxaxis}
\end{tikzpicture}
\[
 z_m<\alpha
\]
\end{minipage}
\hfill
\begin{minipage}{0.48\textwidth}
\centering
\begin{tikzpicture}
% Same qualitative plot, but now the maximum of g_2 occurs before z_m.
% The numerical values are rescaled only to keep both panels visually comparable.
\def\alphaval{0.035}
\def\Cmplot{0.8}
\def\zone{0.05311}
\begin{semilogxaxis}[
    width=\textwidth,
    height=0.72\textwidth,
    axis lines=left,
    xlabel={},
    ylabel={value},
    xmin=0.03,
    xmax=20,
    ymin=0,
    ymax=3.0,
    samples=400,
    domain=0.03:20,
    legend style={draw=none, fill=none, at={(0.98,0.98)}, anchor=north east, font=\scriptsize},
    xtick={\alphaval,\zone},
    xticklabels={$\alpha$,$z_m$},
    ytick=\empty,
    tick style={draw=none},
    clip=false,
]
  \addplot[black, thick]
    {(1/6 + x/3)/(\alphaval + x)};
  \addlegendentry{$g_1(z)=\frac{1/6+z/3}{\alpha+z}$}

  \addplot[black, thick, dashed]
    {\Cmplot*sqrt(x)/(\alphaval + x)};
  \addlegendentry{$g_2(z)=\frac{C_m\sqrt z}{\alpha+z}$}

  \draw[densely dotted] (axis cs:\zone,0) --
    (axis cs:\zone,{(1/6+\zone/3)/(\alphaval+\zone)});
  \draw[densely dotted] (axis cs:\alphaval,0) --
    (axis cs:\alphaval,{\Cmplot*sqrt(\alphaval)/(\alphaval+\alphaval)});
\end{semilogxaxis}
\end{tikzpicture}
\[
 \alpha<z_m
\]
\end{minipage}
\caption{Comparison of \(g_1\) and \(g_2\).}
\label{fig:g1-g2-comparison}
\end{figure}

It remains to compare (16) with the old bound.  If
\(z_m\le\alpha<1/2\), then \(g_1\) is strictly decreasing, so
\[
 \frac{1+2z_m}{6(\alpha+z_m)}
 =g_1(z_m)<g_1(0)=\frac{1}{6\alpha}.
\]
If \(0<\alpha\le z_m\), equation (13) gives
\[
 3C_m\sqrt{z_m}=\frac12+z_m<1.
\]
Here \(z_m<1/2\) follows immediately from \(C_m\ge2\) and (14).
Therefore
\[
 3C_m\sqrt\alpha\le3C_m\sqrt{z_m}<1,
\]
which is equivalent to
\[
 \frac{C_m}{2\sqrt\alpha}<\frac{1}{6\alpha}.
\]
This proves (17).

Finally, substituting \(\alpha=1/2\) into the second branch of (16) gives
\[
 \frac{1+2z_m}{6(1/2+z_m)}=\frac13,
\]
and hence (18).
\end{proof}
\end{document}
